\section{Introducción al Dossier Comercial y de Cumplimiento}

Este dossier presenta la documentación comercial, el manual operativo de campo y el marco de cumplimiento regulatorio del collar de geocercamiento ganadero DR-VF-01. Incluye la hoja de datos ejecutiva (\textit{Datasheet}), el manual de instalación y configuración, las certificaciones de compatibilidad electromagnética (EMC), seguridad eléctrica y las normativas ambientales para su escalamiento comercial.

\section{Datasheet Ejecutivo del Producto}

\subsection{Descripción General}

El DR-VF-01 es una solución de grado industrial para el geocercamiento y monitoreo continuo de ganado extensivo. Integra geolocalización multiconstelación GNSS, conectividad LoRa Sub-GHz de ultra-largo alcance, análisis de movimiento mediante IMU y una etapa disuasiva progresiva controlada en tiempo real.

\subsection{Características Principales}

\subsubsection{Geolocalización Multiconstelación de Alta Sensibilidad}

\begin{itemize}[leftmargin=*]
    \item \textbf{Receptor GNSS}: u-blox MAX-M10S con recepción simultánea GPS, GLONASS, Galileo y BeiDou.
    \item \textbf{Precisión Posicional}: $< 3.0\,\text{m}$ RMS en condiciones de vista despejada.
    \item \textbf{Tiempo de Adquisición (TTFF)}: $< 30\,\text{s}$ en arranque en frío (\textit{Cold Start}) y $< 2\,\text{s}$ en \textit{Hot Start}.
\end{itemize}

\subsubsection{Enlace de Comunicación LoRa Sub-GHz}

\begin{itemize}[leftmargin=*]
    \item \textbf{Transceptor RF}: Semtech SX1262 para red LoRaWAN o P2P.
    \item \textbf{Alcance Operativo}: Hasta $15.0\,\text{km}$ en línea de vista (\textit{LoS}) en entornos rurales.
    \item \textbf{Banda de Frecuencia y Potencia}: $915\,\text{MHz}$ (banda ISM) con potencia de salida programable hasta $+22\,\text{dBm}$.
\end{itemize}

\subsubsection{Detección Inteligente de Movimiento y Wake-Up}

\begin{itemize}[leftmargin=*]
    \item \textbf{Sensor Inercial}: IMU LSM6DSOX de 6 ejes con procesamiento integrado de eventos.
    \item \textbf{Activación Selectiva}: Interrupción por hardware cuando el bovino supera el umbral inercial.
    \item \textbf{Umbral Configurable}: Ajustable en rango de $0.1\,g$ a $1.0\,g$.
\end{itemize}

\subsubsection{Autonomía y Gestión Solar de Energía}

\begin{itemize}[leftmargin=*]
    \item \textbf{Autonomía Base}: Hasta $80\,\text{días}$ en batería LiPo ($1500\,\text{mAh}$) sin recarga solar.
    \item \textbf{Consumo en Reposo}: $< 5.0\,\mu\text{A}$ en modo \textit{Deep Sleep}.
    \item \textbf{Recarga Solar MPPT}: Cargador dedicado CN3791 acoplado a panel monocristalino de $2.0\,\text{W}$.
\end{itemize}

\subsection{Especificaciones Técnicas del Sistema}

\begin{table}[H]
\centering
\caption{Matriz de componentes principales de la plataforma DR-VF-01}
\label{tab:componentes_principales}
\vspace{2mm}
\begin{tabularx}{\textwidth}{X X}
\toprule
\rowcolor{categorybg}
\textbf{Módulo / Subgrupo} & \textbf{Especificación Técnica} \\
\midrule
Microcontrolador Core & ESP32-S3-WROOM-1U-N16R8 ($16\,\text{MB}$ Flash, $8\,\text{MB}$ PSRAM) \\
\rowcolor{rowalt}
Navegación GNSS & u-blox MAX-M10S (Multiconstelación 72 canales) \\
Radiofrecuencia LoRa & Semtech SX1262 ($915\,\text{MHz}$, $+22\,\text{dBm}$) \\
\rowcolor{rowalt}
Sensor de Movimiento & STMicroelectronics LSM6DSOX IMU 6 ejes \\
Gestión Fotovoltaica & CN3791 MPPT Solar Charger ($6\,\text{V}$) \\
\rowcolor{rowalt}
Batería de Respaldo & LiPo $3.7\,\text{V}$ / $1500\,\text{mAh}$ (Protección PCM integrada) \\
Elemento Radiante & Antena FPC interna omnidireccional ($+2.5\,\text{dBi}$) \\
\bottomrule
\end{tabularx}
\end{table}

\begin{table}[H]
\centering
\caption{Parámetros globales de desempeño operativo}
\label{tab:desempeno_sistema}
\vspace{2mm}
\begin{tabularx}{\textwidth}{X r}
\toprule
\rowcolor{categorybg}
\textbf{Parámetro Métrico} & \textbf{Valor / Rango Garantizado} \\
\midrule
Alcance de Enlace RF LoRa & Hasta $15.0\,\text{km}$ (Línea de Vista en Campo) \\
\rowcolor{rowalt}
Precisión Posicional GNSS & $< 3.0\,\text{m}$ RMS \\
Tiempo de Respuesta Algoritmo PiP & $< 50\,\text{ms}$ \\
\rowcolor{rowalt}
Consumo de Corriente en \textit{Deep Sleep} & $< 5.0\,\mu\text{A}$ \\
Rango de Temperatura Operativa & $-20\,^\circ\text{C}$ a $+65\,^\circ\text{C}$ \\
\rowcolor{rowalt}
Grado de Estanqueidad Estructural & IP67 (IEC 60529) \\
\bottomrule
\end{tabularx}
\end{table}

\section{Manual de Operación e Instalación en Campo}

\subsection{Procedimiento de Ensamble y Ajuste}

\subsubsection{Fase 1: Preparación y Carga Inicial}

\begin{enumerate}[leftmargin=*]
    \item Retire el tapón protector del conector USB-C e inserte el cable de carga suministrado.
    \item Conecte la fuente de alimentación externa ($5\,\text{V DC} / 1\,\text{A}$).
    \item Monitoree el LED indicador de estado hasta observar color verde continuo (duración estimada de carga total: $2.0\,\text{horas}$).
    \item Reinstale cuidadosamente el tapón de silicona asegurando el sellado IP67.
\end{enumerate}

\subsubsection{Fase 2: Fijación y Ergonomía sobre el Bovino}

\begin{enumerate}[leftmargin=*]
    \item Posicione la correa de poliuretano alrededor de la sección superior del cuello del animal.
    \item Verifique que el gabinete electrónico principal quede orientado en la cresta superior del cuello para optimizar la captación solar y la cobertura GNSS/RF.
    \item Ajuste la hebilla de seguridad dejando una holgura equivalente al paso de una mano ($2\text{--}3\,\text{cm}$) entre la correa y la piel para prevenir laceraciones.
    \item Asegure el retenedor secundario y recorte el remanente de correa sobrante.
\end{enumerate}

\subsection{Configuración de Geocercas vía Servidor Web Embebido}

\subsubsection{Acceso al Servidor Local SoftAP}

\begin{enumerate}[leftmargin=*]
    \item Encienda la interfaz Wi-Fi en el dispositivo móvil de configuración.
    \item Seleccione el punto de acceso denominado \texttt{DR-VF-01-CONFIG}.
    \item Abra un navegador e ingrese a la dirección IP \url{http://192.168.4.1}.
    \item Autentíquese con las credenciales por defecto: usuario \texttt{admin}, contraseña \texttt{geofence2026}.
\end{enumerate}

\subsection{Codificación de Indicadores LED de Diagnóstico}

\begin{description}[leftmargin=*, style=nextline, font=\bfseries\ttfamily]
    \item[Verde Fijo] Proceso de carga de batería activo vía USB-C / Solar.
    \item[Verde Parpadeante] Batería completamente cargada ($V_{\text{bat}} \ge 4.18\,\text{V}$).
    \item[Azul Parpadeante] Modo de configuración SoftAP y servidor HTTP activo.
    \item[Rojo Fijo] Alerta de batería crítica ($V_{\text{bat}} \le 3.40\,\text{V}$).
    \item[Rojo Parpadeante] Evento de rebasamiento de geocerca y activación de etapa disuasiva.
    \item[Amarillo Parpadeante] Excepción de sistema o pérdida de señal GNSS (\textit{Timeout}).
\end{description}

\section{Normativas de Certificación y Ensayos de Cumplimiento}

\subsection{Certificaciones de Radiofrecuencia (RF)}

\subsubsection{Regulación FCC (Part 15 Subpart C)}

El collar DR-VF-01 cumple con los límites de emisión radiada para dispositivos intencionales en la banda ISM de $902\text{--}928\,\text{MHz}$ bajo FCC Part 15.247.

\subsubsection{Directiva Europea RED (2014/53/EU)}

Evaluación de conformidad según la directiva RED para equipos de radio, garantizando un uso eficiente del espectro electromagnético en frecuencias sub-GHz.

\subsection{Seguridad Eléctrica y Disuasión}

\subsubsection{Cumplimiento IEC 62368-1 / IEC 60950-1}

Cumplimiento de requisitos de seguridad en equipos de tecnología de la información, abarcando protección contra choques eléctricos, sobrecalentamiento e inflamabilidad.

\subsubsection{Norma de Impulsos Disuasivos IEC 60335-2-76 (Anexo CC)}

La etapa de disuasión del collar cumple los límites de energía pico ($< 0.2\,\text{J}$) y tensión máxima ($< 2.5\,\text{kV}$) establecidos para cercas eléctricas de control animal, garantizando la seguridad biológica del bovino sin causar quemaduras ni secuelas fisiológicas.

\subsection{Compatibilidad Electromagnética (EMC)}

\subsubsection{Ensayos de Emisión}

\begin{itemize}[leftmargin=*]
    \item \textbf{FCC Part 15 Subpart B}: Dispositivo digital Clase B.
    \item \textbf{EN 55032 / EN 55035}: Emisiones e inmunidad en equipos multimediales.
\end{itemize}

\subsubsection{Ensayos de Inmunidad}

\begin{itemize}[leftmargin=*]
    \item \textbf{IEC 61000-4-2}: Inmunidad a descargas electrostáticas (ESD) hasta $\pm 8\,\text{kV}$ por contacto y $\pm 15\,\text{kV}$ por aire.
    \item \textbf{IEC 61000-4-3}: Inmunidad a campos electromagnéticos radiados ($10\,\text{V/m}$).
    \item \textbf{IEC 61000-4-4}: Inmunidad a ráfagas de transitorios rápidos (\textit{EFT/Burst}).
    \item \textbf{IEC 61000-4-5}: Inmunidad a sobretensiones conducidas (\textit{Surge}).
\end{itemize}

\section{Cumplimiento Ambiental y Logística}

\subsection{Directiva RoHS y REACH}

El dispositivo cumple con la Directiva 2011/65/EU (RoHS 2/3) y la Regulación REACH, certificando que todos los componentes y la PCB están libres de plomo, mercurio, cadmio y retardantes de llama bromados.

\subsection{Normativa de Transporte de Baterías de Litio}

\subsubsection{Manual de Pruebas UN 38.3}

La celda LiPo integrada superó las pruebas de transporte de la ONU (altitud, choque térmico, vibración, impacto, sobrecarga y cortocircuito externo).

\subsubsection{Regulación IATA / ICAO}

Empaque y etiquetado conforme a la instrucción de embalaje PI967 (Baterías de ion de litio contenidas en equipos).

\section{Garantía, Soporte Técnico y Mantenimiento}

\subsection{Esquema de Garantía de Fábrica}

\begin{itemize}[leftmargin=*]
    \item \textbf{Dispositivo Electrónico y Carcasa}: 24 meses contra defectos de fabricación o fallas de estanqueidad.
    \item \textbf{Batería LiPo Integrada}: 12 meses o 500 ciclos completos de carga/descarga.
    \item \textbf{Soporte Técnico Especializado}: 12 meses de soporte remoto y diagnósticos de red sin costo.
\end{itemize}

\subsection{Política de Actualizaciones Firmware OTA}

\begin{itemize}[leftmargin=*]
    \item Soporte de actualización remota Over-The-Air (OTA) cifrada vía LoRaWAN / SoftAP.
    \item Repositorio de parches de seguridad y optimización de algoritmos PiP.
\end{itemize}

\section{Conclusiones}

El Dossier Comercial y de Cumplimiento consolida la viabilidad técnica y normativa del collar DR-VF-01. Las certificaciones de seguridad eléctrica (IEC 62368-1 e IEC 60335-2-76) y la inmunidad EMC garantizan la seguridad operativa en campo abierto. Asimismo, la hoja de datos técnicos y el manual operativo aseguran un despliegue ágil, confiable y escalable para el sector agropecuario.